Read in Gauge Theory Gravity through Fact \claimref{synthesis.imported-gtg-dictionary}(i): the coframe is the physical-vector map $\hb^{-1}$ (in components $e^a(v) = \gamma^a\cdot\hb^{-1}(v)$), the derived metric is $g(a,b) = \hb^{-1}(a)\cdot\hb^{-1}(b)$, and the dual frame is the background directions $E_a = \hb(\gamma_a)$. (i) Write $r$ for a tangent vector with $k(r) = 1$. By the residue, $G^{-1}(r, r) = \tau^{-1}/T + O(1)$, so $Q(\hb^{-1}(r)) = \tau^{-1}/T + O(1)$, and at least one component $\gamma^a\cdot\hb^{-1}(r)$ of the physical vector $\hb^{-1}(r)$ is of order $\tau^{-1/2}$. The dual frame $E_a = \hb(\gamma_a)$ then has $k(E_a) = O(\tau^{1/2})$ for that $a$: its image collapses into $\ker k$ at $\tau = 0$. A soldering form with a radical would have the derived metric $g$ finite and degenerate; here $g(r,r) = Q(\hb^{-1}(r))$ is infinite. (ii) In the weighted module the Gram of the coframe is the matrix $G(\tau)$ of \cite[proof of Thm.~6.2]{ButlerT4} (\claimref{sheaf.metrics-residue-square-root-twist}), $\begin{psmallmatrix}\lambda & b^{\mathsf T}\\ b & G_1\end{psmallmatrix}$ with $\lambda(0) = b(0) = 0$ and $G_1(0)$ invertible (by \claimref{sheaf.tameness}): that is, the internal cometric has Gram matrix $\eta(0) = \mathrm{diag}(0, \eta_3)$. For the stabiliser write $A = \begin{psmallmatrix}\alpha & u^{\mathsf T}\\ v & B\end{psmallmatrix}$; then $A\eta + \eta A^{\mathsf T} = \begin{psmallmatrix} 0 & u^{\mathsf T}\eta_3\\ \eta_3 u & B\eta_3 + \eta_3 B^{\mathsf T}\end{psmallmatrix}$, which vanishes if and only if $u = 0$ and $B \in \mathfrak{so}(\eta_3)$, with $\alpha$ and $v$ free: $1 + 3 + 3 = 7$. (iii) Take the ambient basis $(k, m_1, m_2, m_3, k^*, w)$, $k^*$ the partner of the radical covector over $\Sigma$ (\claimref{sheaf.localised-transport}), with $\ip{k}{k^*} = 1$, $\ip{m_i}{m_j} = (\eta_3)_{ij}$, $Q(w) = +1$ and all other products zero, which has signature $(2,4)$ for $\eta_3$ negative definite. An element of $\mathfrak{so}(2,4)$ preserving $\mathcal{P} = \mathrm{span}(k, m_i)$ preserves $\mathcal{P}^\perp = \mathrm{span}(k, w)$; solving the linear conditions gives an eight-dimensional algebra whose restriction to $\mathcal{P}$ is the seven-dimensional algebra of (ii) (the one-dimensional kernel acts on $\mathrm{span}(k^*, w)$ alone). The element restricting to $\mathrm{diag}(\alpha, 0, 0, 0)$ on $\mathcal{P}$ is forced by $\ip{k}{k^*} = 1$ to send $k^* \mapsto -\alpha k^*$ and, by orthogonality, to fix $w$: it is the boost of the hyperbolic plane $\mathrm{span}(k, k^*)$. The script verifies both dimensions, the surjectivity of the restriction, and the lift of the dilation as exact matrices. (iv) The first case is Theorem \claimref{dark-energy.walls-free-sector}(a) on the invertible region (rank $24$ and $16$); the second is Theorem \claimref{dark-energy.walls-free-sector}; the third is (i) together with the identification $e = \nabla X$ of Lemma \claimref{dark-energy.stelle-west-vector-frozen-point}, whose Gram in the tractor chart is the derived metric $g(a,b) = \hb^{-1}(a)\cdot\hb^{-1}(b)$, which has the pole.
